\documentclass[../../main.tex]{subfiles}
\begin{document}

{\bf [解答]}

与えられた関数
\begin{align}
  F(\theta)=\int_0^{\theta}x\cos(x+\alpha)dx \quad  (0<\alpha<\dfrac{\pi}{2},0\le \theta \le \dfrac{\pi}{2}) \label{eq:1}
\end{align}
について考える．一階微分は
\begin{align}
  F'(\theta)=\theta\cos(\theta+\alpha)
\end{align}
だから$F(\theta)$の増減表は以下のようになる．

\begin{table}[htb]
  \centering
  \caption{$F(\theta)$の増減表}
  \begin{tabular}{c|ccccc}
    $\theta$ & $0$ & & $\dfrac{\pi}{2}-\alpha$ & & $\pi/2$ \\
    \hline
    $F'$ & & $+$ & $0$ & $-$ & \\
    \hline
    $F$ & & $\nearrow$ & & $\searrow$ & \\
  \end{tabular}
\end{table}

したがって，$F(\theta)$は$\theta=\dfrac{\pi}{2}-\alpha$で最大である．
$F(\theta)$の積分を計算すると
\begin{align}
  F(\theta)
  &=\Bigl[x\sin(x+\alpha)+\cos(x+\alpha)\Bigr]_0^{\theta} \\
  &=\theta\sin(\theta+\alpha)+\cos(\theta+\alpha) - \cos\alpha
\end{align}
だから，$\theta=\dfrac{\pi}{2}-\alpha$を代入して
この時の値は
\begin{align}
  \max F(\theta)
  &= F(\dfrac{\pi}{2}-\alpha) \\
  &=\left(\dfrac{\pi}{2}-\alpha\right)-\cos\alpha
\end{align}
である．

\newpage

\end{document}
